\section{Introduction}
A spectrum analyzer need not generate audio to interfere with audio
processing. If its FFT runs on the audio engine thread, the engine must finish
that work before it can deliver the next block of sound. Computing a spectrum
only a few dozen times per second can make the average cost look modest while
leaving occasional blocks much more expensive than their neighbors. Several
analyzers updating together can compound the burst.

This is the problem that motivated Fourier, a spectrum analyzer for the VCV
Rack modular synthesizer. Rack invokes a module's processing function once
per sample~\cite{rack}. The FFT was therefore built to advance a little on
each invocation: execute some butterflies, save the current position, and
resume when the next sample arrives. Analysis remains in the ordinary
processing loop. The display receives the completed spectrum later, while
audio processing continues between pieces of the transform.

The distinction is between the amount of work and when that work happens.
A conventional radix-2 FFT takes $O(N\log N)$ operations for $N$ samples.
Making it resumable does not reduce that count. It makes the computation
available for scheduling over an interval in which the application can afford
to wait for the result. Spectral visualization is a natural application:
its output is a sequence of display updates rather than an audio sample that
must be returned immediately. Whether the added delay is acceptable still
depends on the desired display response.

Scheduling butterflies alone is insufficient. Windowing and arranging the
input, reconstructing a real spectrum, computing magnitudes, and preparing
output can each leave a large loop at a frame boundary. The implementation
studied here distributes those operations too. It selects a frame, preserves
its input while new samples arrive, advances the analysis in dependency order,
and publishes the completed result at the end of a fixed number of calls.
Figure~\ref{fig:schedule} illustrates this change in execution.

There is a competing advantage to doing the work at once: optimized FFT
libraries use vectorization and carefully organized memory access. A custom
resumable transform may give up some of that efficiency. A useful evaluation
must therefore ask both whether distribution reduces bursts within the same
implementation and whether the resulting tradeoff remains worthwhile against
native libraries. Comparing only a custom batch FFT with its incremental
version would answer the first question but leave the second open.

This paper makes three contributions. It presents an implementable schedule
for complete windowed analysis within a sample-processing loop; explains its
work bound, input lifetime, and publication delay; and measures the tradeoff
in controlled analysis and complete Rack modules. The method uses established
FFT mathematics and cooperative scheduling. Its contribution is their
application to the complete analyzer and the evidence needed to judge that
choice, rather than a new Fourier transform or a claim that distribution is
always faster.

The results support a focused conclusion: distributing analysis can buy
substantial processing headroom in burst-heavy configurations, but that
headroom has costs in total work and result age. Native batch and hybrid
implementations remain useful alternatives. Host scheduling also matters:
lower typical processing peaks did not consistently produce fewer late blocks
in the measured system.
